
\textbf{Classical PK and NLME Modelling}. Population PK is typically formulated with nonlinear mixed-effects (NLME) models, where a compartmental ODE is shared across subjects and parameters decompose into fixed (population) and random (individual) effects~\citep{lavielle2014}. Recent ML extensions include GP and neural mixed-effects formulations, including mixed-effects neural ODEs \cite{kipnis2025metabeta}, but these approaches are often developed for data-rich settings rather than the small, sparse studies common in early drug development. And different from us, focus on learning parameters of mechanistic models \cite{haggstrom2025simulation}.

\textbf{Amortized Inference for PK}. Simulation-based amortization has recently been used to accelerate pharmacokinetic inference by pretraining neural estimators on simulated data~\citep{pmlr-v235-arruda24a}. In Arruda et al., amortization is primarily performed in \emph{parameter space}: the learned estimator approximates posterior inference over the scalar parameters of a specified PK simulator, which are then propagated through the mechanistic model to obtain concentration--time predictions. This is complementary to the notion of amortization pursued in \texttt{AICMET}, which maps study context directly to posterior-predictive concentration functions in a zero-shot manner~\citep{marin2025amortized}, but retains a neural-process-style latent-variable representation of study and individual effects. Building on this line, \texttt{PFF} amortizes the posterior-predictive operator directly in function space. It learns a conditional flow over continuous-time concentration functions, conditioned on both study-level context and individual partial observations, and therefore avoids dataset-specific retraining while targeting the posterior-predictive distribution itself rather than first inferring simulator parameters.


\textbf{Generative Modelling for Continuous Functions}. Function-space generative models learn continuous-time transport maps between distributions over trajectories, including diffusion-based and flow-matching approaches~\citep{bilovs2023modeling,kerrigan2023functional,kollovieh2024flow,jahn2025trajectory}. Our method builds most directly on the dynamic conditional optimal-transport perspective for functional flow matching~\citep{kerrigan2024dynamic}.

%\textbf{Classical PK and NLME Modelling}. Classical population pharmacokinetics is built around nonlinear mixed-effects (NLME) models, where a mechanistic compartmental ODE is shared across subjects while parameters decompose into fixed (population-level) and random (individual-level) effects. 
%
%\citet{lavielle2014} is the standard reference for this family of approaches.
%
%Beyond the classical toolbox, machine-learning variants span (i) nonparametric Gaussian-process formulations of mixed-effects regression~\citep{shi2012mixed}; (ii) hybrid models that keep the mixed-effects decomposition, but model shared structure with neural networks while using GPs for individual deviations~\citep{chung2020deep}; and 
%
%(iii) mixed-effects neural ODE approaches~\citep{nazarovs2022mixed}.
%
%In pharmacology-specific work, neural PK/PD models show that neural ODE architectures can learn latent dynamics and predict response time courses, but they are typically developed in data-rich regimes, and thus are \textit{not} tailored to the small, sparse studies common in early drug development.

%\textbf{Amortized Inference for PK}. A different axis of work targets the cost of inference itself. 
%
%\citet{pmlr-v235-arruda24a} propose an amortized simulation-based pipeline for NLME.
%
%In their work a neural posterior estimator is pretrained on simulations to produce individual-level posteriors, and a separate population-level likelihood approximation is then built on top to infer study-level parameters for a given dataset.
%
%Amortization therefore accelerates per-subject inference, but a study-level fit is still performed anew per dataset.
%
%\texttt{AICMET} pushes amortization further by learning to map the study context directly to posterior-predictive distributions over concentration functions in a \textit{zero-shot} fashion, in the spirit of prior-fitted networks~\citep{marin2025amortized}.
%
%Yet \texttt{AICMET} inherits limitations of neural-process-style latent variable modeling.

%\textbf{Generative Modelling for Continuous Functions}. A rapidly growing literature studies generative modeling directly in function space, where one learns a continuous-time transport map from a simple reference measure over functions to a target distribution.
%
%Recent directions include diffusion-style constructions for temporal data treated as continuous functions~\citep{bilovs2023modeling}, functional flow-matching formulations~\citep{kerrigan2023functional, kollovieh2024flow}, and alternative stochastic dynamics (\textit{i.e.}, jump-diffusion constructions) for matching trajectory distributions~\citep{jahn2025trajectory}.
%
%Our work builds most directly on the dynamic conditional OT viewpoint for functional flow matching introduced by~\citet{kerrigan2024dynamic}.

% Generative Models for Time Series \cite{kollovieh2024flow,bilovs2023modeling,jahn2025trajectory}

% Machine learning solutions for NLME modelling span a wide range of methodologies, from traditional nonparametric Gaussian process regression methods \citep{shi2012mixed} to more scalable variants that condition the GP mean functions on neural networks \citep{chung2020deep}. 
% %
% Neural ODE-based approaches have also been applied in this context \citep{nazarovs2022mixed}, although they typically impose a linear dependence on the fixed effects, and still rely on computationally expensive and often unstable adjoint methods for training. 
% %
% Specifically tailored solutions for pharmacokinetics \citep{lu2021deep} have demonstrated that neural ODEs can learn both complex latent dynamics and response times as a function of dosing schemes. 
% %
% However, these methods are generally trained on large datasets and ignore the hierarchical, population-level structure central to NLME modeling.

% Our proposal provides an alternative to these approaches by combining compartmental priors with amortized in-context inference, enabling zero-shot adaptation from sparse, irregularly sampled data while explicitly capturing both population- and individual-level variability.

% Transformer-based models, on the other hand, have recently been explored for their capacity to perform Bayesian inference via in-context learning \citep{muller2021transformers, mittal2025amortized}. In particular, \citet{muller2021transformers} introduce the notion of prior-fitted networks, drawing a connection to simulation-based inference where domain-specific priors are incorporated, and generalization is achieved through zero-shot adaptation to sparse data. In contrast, \citet{mittal2025amortized, mittal2025context} present a comprehensive evaluation across training objectives and architectures, highlighting the amortized nature of these methods—allowing practitioners to bypass costly retraining when faced with new datasets. Neural Processes \citep{garnelo2018neural}, which emphasize inference over unstructured latent variables, have also benefited from the integration of attention mechanisms and transformer architectures \citep{kim2019attentive, nguyen2022transformer}.

% In the context of dynamical systems, pretraining has proven effective for both time-series forecasting \citep{dooley2024forecastpfn} and the recovery of governing equations in physical systems \citep{dascoli2024odeformer}. These models are typically pretrained on large synthetic datasets, where the practitioner defines the family of processes to be recovered. Some approaches aim directly at predictive distributions \citep{ansari2024chronos}, while others attempt to infer the underlying generators. The latter has shown promise in modeling jump diffusion processes \citep{berghaus2024foundation} and stochastic differential equations \citep{seifner2025foundation}.

% Our approach builds on these insights. Similar to \textit{prior-fitted networks}, we restrict the class of inferred systems by introducing a simulation-based prior that emphasizes compartment PK models, and places a stationary Ornstein–Uhlenbeck (OU) process on the parameters of the PK model. To account for the hierarchical structure of mixed-effects models, we incorporate unstructured latent variables for both fixed and random effects—drawing inspiration from neural processes—while retaining the flexibility of neural representations by not enforcing permutation invariance. In contrast to standard inference methods, which often assume fully observed trajectories, we focus on the posterior predictive distribution of partially observed systems. This formulation enables both, the generation of trajectories for unseen individuals and the forecasting for partially observed ones. Consequently, our model leverages the learned fixed and random effect representations in a unified generative and predictive framework.